%======================================================================
% Section 10: Energy Transport and Dissipation
%======================================================================
\section{Energy Transport and Dissipation}
\label{sec:energy}

Constitutive theories carry energy-storage, transport, and dissipation consequences, but the expressions depend on their state variables and field equations. This section derives the classical local flux under explicit conventions, states the interface balance at a general level, and records family-specific flux claims as open unless their free energy, constitutive evolution, boundary power, and dissipation inequality have been checked. No universal generalized-continuum flux is assumed.

\subsection{General energy balance}
\label{sec:10.1}

For classical thermoelasticity, with specific internal energy $e$, the total energy per unit volume is

\begin{equation}
\mathcal{E} = \frac{1}{2}\rho v_i v_i + \rho e,
\label{eq:10.1}
\end{equation}

and, with no internal heat source, the exact Eulerian balance is

\begin{equation}
\frac{\partial\mathcal{E}}{\partial t} + \partial_i F_i = \rho f_i v_i,
\qquad F_i=q_i-\sigma_{ij}v_j+\mathcal{E}v_i,
\label{eq:10.2}
\end{equation}

where $\rho f_i v_i$ is mechanical power per unit volume and $F_i$ is the exact Eulerian flux for the classical local Cauchy continuum. Equation~\eqref{eq:10.2} is not asserted as a universal generalized-multiphysics identity: additional field energy, external electrical or magnetic power, interface storage, and internal dissipation must be derived for each specified model.

\subsection{Classical energy flux}
\label{sec:10.2}

For first-order small-amplitude fields about a resting reference state, after measuring $e$ relative to a constant $e_0$, the perturbation advection $\rho[\tfrac12 v_jv_j+(e-e_0)]v_i$ is $O(\epsilon^2)$ and is absent from the first-order field balance. The reduced outward-positive flux used below is

\begin{equation}
J_i = q_i - \sigma_{ij}v_j,
\label{eq:10.3}
\end{equation}
It is not the exact Eulerian total-energy flux: the latter is $F_i=J_i+\mathcal{E}v_i$ as in Eq.~\eqref{eq:10.2}. The reduced expression is used only for the stated first-order wave-balance convention; this order argument alone does not justify omitting the advective term from a second-order cycle-averaged total-energy flux. Under the outward-positive convention the mechanical contribution is $-\sigma_{ij}v_j$ and the thermal contribution is $q_i$.

For time-harmonic mechanical fields with convention $e^{-i\omega t}$, the mean of the reduced flux is

\begin{equation}
\langle J_i\rangle_t=\bar q_i-\frac12\Re[\hat{\sigma}_{ij}\hat{v}_j^*],
\label{eq:10.4}
\end{equation}

where $\bar q_i$ is the separately evaluated mean heat flux (zero for a purely oscillatory perturbation about a uniform reference state). The factor $1/2$ and complex conjugate follow from time averaging the real fields. For an interface problem using this reduced convention the quantity is $J_n=J_i n_i$. If a second-order Eulerian total-energy balance is claimed, the mean of the advective contribution must also be evaluated or shown to vanish under an additional model-specific assumption.

A point worth keeping in view is that the reduced classical flux contains \emph{no} term associated with the internal variables that generalized theories introduce. The entire difficulty of Section~\ref{sec:10.4} follows from that absence.





\subsection{Interface energy balance}
\label{sec:10.3}

Consider a stationary interface $\Sigma$ between two media with unit normal $\bm{n}$. Integrating the model-correct total-energy balance over a thin pillbox straddling the interface and taking the thinness limit gives the jump condition

\begin{equation}
\jump{\mathcal{F}_n} = \mathcal{P}_\Sigma - \frac{\partial\mathcal{E}_\Sigma}{\partial t},
\label{eq:10.5}
\end{equation}
where $\mathcal{F}_n=\mathcal{F}_i n_i$ is the normal component of the exact, model-correct Eulerian total-energy flux (equal to $F_n$ for the classical continuum), $\mathcal{P}_\Sigma$ is the mechanical power input at the interface, and $\mathcal{E}_\Sigma$ is the interfacial energy density. The reduced $J_i$ of Eq.~\eqref{eq:10.3} may replace it only within the stated first-order approximation.

For a passive interface with no mechanical power input and no interfacial energy storage,

\begin{equation}
\jump{\mathcal{F}_n} = 0,
\label{eq:10.6}
\end{equation}
that is, the normal model-correct total-energy flux is continuous across the interface. Two qualifications are needed before this condition is used numerically. First, it is a statement about the \emph{correct} energy flux of the theory in which it is used; a flux built from the classical form alone, applied to a family with internal variables, will not satisfy it even when the mathematics is otherwise flawless. Second, the continuity of the normal flux does not by itself imply the scalar conservation identity of Section~\ref{sec:10.5}: that identity is a statement about modal energies, while \eqref{eq:10.6} is a statement about a single component of a flux at a single surface, and the two are related only after the modal decomposition and normalisation conventions have been fixed.

\subsection{Energy transport in generalized theories: the internal-variable structure}
\label{sec:10.4}

The first-law balance retains the boundary power terms for a specified constitutive theory. What changes between generalized theories is the stored free energy and dissipation, which follow from that theory's state variables and evolution laws. There is no universal additional flux term $-\sum_\alpha Y^\alpha\dot z^\alpha$: this has units of power per volume, not energy flux, unless an additional spatial transport structure is specified. A family-specific identity therefore requires the model's free energy, evolution law, and boundary terms.

\subsubsection{Piezothermoelastic energy flux}

For a coupled piezothermoelastic model, the mechanical, electrical, and thermal power terms should be derived from one stated balance convention and the coupled constitutive equations before modal ratios are interpreted.

\subsubsection{Energy flux in generalized theories}

A family-specific energy balance requires the free energy, flux variables, boundary powers, dissipation terms, and constitutive equations of the selected model.

For the surface thermoelastic waves studied by Ivanov and Savova, the energy-flux vector connected with kinetic and internal energy is reported as not parallel to the half-space surface; its direction varies with distance from the surface, its magnitude decays to zero with depth, and it decays exponentially along the direction in which the superposed waves' constant-phase planes intersect the surface (printed p.~564). This is a paper-specific surface-wave result, not a universal flux law \cite{IvanovSavova1991}.

\subsection{Remote-plane energy partition}
\label{sec:10.5}

For passive regions with steady time-averaged powers, no source, no net energy accumulation, and no unlisted energy channel, the remote-plane balance is

\begin{equation}
R_{\mathrm{rem}} + T_{\mathrm{rem}} + A_{\mathrm{bulk},L} + A_{\mathrm{int}} + A_{\mathrm{bulk},R} = 1,
\label{eq:10.11}
\end{equation}

where $R_{\mathrm{rem}}=P_r/P_i$ and $T_{\mathrm{rem}}=P_t/P_i$ use positive incident/reflected/transmitted power magnitudes at explicitly stated remote reference planes. The normalized losses are $A_{\mathrm{bulk},L}=\mathcal{D}_L/P_i$, $A_{\mathrm{int}}=\mathcal{D}_{\Sigma}/P_i$, and $A_{\mathrm{bulk},R}=\mathcal{D}_R/P_i$, with the control-volume powers defined in Eq.~\eqref{eq:10.13}. This is distinct from the interface-plane bookkeeping identity in Section~\ref{sec:9.7}; when bulk attenuation separates the planes, $1-R_{\mathrm{rem}}-T_{\mathrm{rem}}$ includes both bulk losses and the local interface residual. Only when the bulk losses and interface residual vanish does the remote-plane balance reduce to

\begin{equation}
R_{\mathrm{rem}} + T_{\mathrm{rem}} = 1.
\label{eq:10.12}
\end{equation}

Kumar and Kansal report a model-specific interface energy-ratio balance with an interaction term in Eq. (80), supporting inclusion of cross terms rather than a universal identity \cite{KumarKansal2017}. Kumar et al. define ratios using average energy flux per incident-wave type; those definitions specify normalization but do not establish closure \cite{Kumar2024}. A related study describes its eight-model sum as approximately unity for specified imperfect-interface cases; this is model-specific evidence, not an identity \cite{Kumar2024a}.

Equation~\eqref{eq:10.11} is the steady control-volume balance under the stated assumptions; its application requires the theory-specific flux, complete power channels, interface storage/power convention, and common remote reference planes. It is not the interface-plane identity of Section~\ref{sec:9.7}. The lossless reduction additionally requires zero bulk loss, zero interface-local residual, and no net interfacial storage. These conditions are not automatic. In the one-dimensional welded example below, the incident--reflected cross term cancels in the time-averaged normal flux; no zero-flux rule for general coupled or evanescent modes is inferred. Report whether each coefficient is an amplitude or a power ratio and state the flux weighting used.

\subsubsection{Separate interface residual from bulk absorption}

Consider a steady time-averaged one-dimensional control volume and take $+x$ as the positive normal direction. Let $\mathcal{F}_x^{(i)}$, $\mathcal{F}_x^{(r)}$, and $\mathcal{F}_x^{(t)}$ denote the incident-, reflected-, and transmitted-mode contributions to the model-correct Eulerian energy flux at stated remote reference planes, with the exact flux used unless a reduced convention is explicitly justified. Per unit cross-sectional area $A$, define $P_i=A^{-1}\int_{S_i}\langle\mathcal{F}^{(i)}\rangle\cdot\hat{\bm{x}}\,dS>0$, $P_r=-A^{-1}\int_{S_r}\langle\mathcal{F}^{(r)}\rangle\cdot\hat{\bm{x}}\,dS>0$, and $P_t=A^{-1}\int_{S_t}\langle\mathcal{F}^{(t)}\rangle\cdot\hat{\bm{x}}\,dS>0$. The selected remote-plane decomposition must make the net left-plane flux equal to $P_i-P_r$; if an incident--reflected cross term remains there, it must be included explicitly in the power balance. Let $\mathcal{F}_L$ and $\mathcal{F}_R$ be the net total-field fluxes on the left and right faces of the interface, both oriented in $+x$. For passive, source-free regions with no bulk or interfacial mean energy accumulation, the bulk and interfacial dissipated powers per unit area are

\begin{equation}
\mathcal{D}_L=P_i-P_r-\mathcal{F}_L,\qquad
\mathcal{D}_{\Sigma}=\mathcal{F}_L-\mathcal{F}_R,\qquad
\mathcal{D}_R=\mathcal{F}_R-P_t.
\label{eq:10.13}
\end{equation}

Summing the three control-volume balances gives

\begin{equation}
P_i=P_r+P_t+\mathcal{D}_L+\mathcal{D}_{\Sigma}+\mathcal{D}_R.
\label{eq:10.14}
\end{equation}

Thus, with $R_{\mathrm{rem}}=P_r/P_i$ and $T_{\mathrm{rem}}=P_t/P_i$, the remote-plane residual $1-R_{\mathrm{rem}}-T_{\mathrm{rem}}$ already contains both bulk dissipation and local interface dissipation. By contrast, $A_{\mathrm{int}}=\mathcal{D}_{\Sigma}/P_i=(\mathcal{F}_L-\mathcal{F}_R)/P_i$ is local to the interface. Adding the bulk terms to $1-R_{\mathrm{rem}}-T_{\mathrm{rem}}$ a second time double-counts them; this follows directly by summing Eqs.~\eqref{eq:10.13}--\eqref{eq:10.14}. The derivation assumes steady averaged powers, no storage in the bulk control volumes, and no unlisted energy channel; other theories require their complete energy flux and power terms.

\subsubsection{Worked one-dimensional welded-interface system}

For a normal-incidence longitudinal wave at $x=0$ in two lossless one-dimensional bars, let $c_m$ be each bar's longitudinal wave speed and take real positive impedances $Z_m=\rho_m c_m$ and wave numbers $k_m=\omega/c_m$. Under the $e^{-i\omega t}$ convention, the displacement amplitudes are

\begin{equation}
\hat u_1=A_i e^{ik_1x}+A_r e^{-ik_1x},\qquad
\hat u_2=A_t e^{ik_2x}.
\label{eq:10.15}
\end{equation}

Displacement and traction continuity, using $\hat\sigma_{xx}=\rho_m c_m^2\partial_x\hat u_m$, give $A_i+A_r=A_t$ and $Z_1(A_i-A_r)=Z_2A_t$. With $A_r$ and $A_t$ as the two unknown amplitudes, these are two scalar conditions and yield the $2\times2$ system

\begin{equation}
\begin{bmatrix}1&-1\\-Z_1&-Z_2\end{bmatrix}
\begin{bmatrix}A_r\\A_t\end{bmatrix}
=-A_i\begin{bmatrix}1\\Z_1\end{bmatrix}.
\label{eq:10.16}
\end{equation}

Its determinant is $-(Z_1+Z_2)$, so the stated positive-impedance case has a unique solution, with amplitude ratios $r=A_r/A_i=(Z_1-Z_2)/(Z_1+Z_2)$ and $t=A_t/A_i=2Z_1/(Z_1+Z_2)$. For the mechanical flux $\langle J_x\rangle=-\tfrac12\Re[\hat\sigma_{xx}\hat v_x^*]$, the incident and reflected power ratios are $R=|r|^2$ and $T=(Z_2/Z_1)|t|^2$. These coincide with the remote-plane ratios because this welded example has lossless, non-attenuating half-spaces. Direct substitution gives $R+T=1$; this result is not asserted for an imperfect, dissipative, or generalized interface. To show the incident/reflected cross term rather than assume it away, the same-medium fields in Eq.~\eqref{eq:10.15} give $\hat v_i=-i\omega A_i e^{ik_1x}$, $\hat\sigma_i=ik_1\rho_1c_1^2 A_i e^{ik_1x}$, $\hat v_r=-i\omega A_r e^{-ik_1x}$, and $\hat\sigma_r=-ik_1\rho_1c_1^2 A_r e^{-ik_1x}$. Before taking the real part in the harmonic flux, their mixed product is
\begin{equation*}
\hat\sigma_i\hat v_r^*+\hat\sigma_r\hat v_i^*=\rho_1c_1^2 k_1\omega\left[-A_iA_r^*e^{2ik_1x}+A_rA_i^*e^{-2ik_1x}\right],
\end{equation*}
the negative of a complex number plus its conjugate, and therefore purely imaginary; hence its time-averaged normal-flux contribution $-\tfrac12\Re(\hat\sigma_i\hat v_r^*+\hat\sigma_r\hat v_i^*)$ is zero. This cancellation is derived only for the stated one-dimensional, lossless longitudinal waves. Direct substitution verifies the matrix dimensions, determinant, amplitudes, field-derived flux weights, cross term, and partition identity for this specified one-dimensional case.

\subsection{Dissipation mechanisms}
\label{sec:10.6}

Dissipative mechanisms can include heat-flux relaxation, viscous stress, dielectric loss, interfacial friction, or hereditary memory. Their presence and non-negative dissipation conditions depend on the precise constitutive equations and admissible state space. A single scalar dissipation function or common parameter-ordering rule does not apply uniformly across these families. In particular, the condition $\tau_q\geq\tau_T\geq0$ and a positive Lord--Shulman relaxation time are not, by themselves, a full proof of thermodynamic admissibility or stability. Deriving model-specific free energies and dissipation inequalities remains open.

\subsection{Thermodynamic admissibility: well-posedness, causality, positivity}
\label{sec:10.7}

Thermodynamic admissibility, mathematical well-posedness, and finite-speed propagation are distinct questions. A constitutive theory requires its own free-energy/entropy statement and dissipation inequality; an initial-boundary-value problem requires a separate existence, uniqueness, and continuous-dependence analysis; propagation speed is a property of the selected equations and parameter regime. Fourier heat conduction gives a parabolic equation with instantaneous mathematical support. Statements for relaxation, phase-lag, memory, fractional, and gradient models depend on their exact constitutive laws.

\subsection{Comparative table: energy transport across theory families}

Table~\ref{tab:energy_transport} summarizes the classical flux and model-specific terms relevant to generalized energy balances.



\begin{table}[htbp]
\centering
\caption{Energy-flux structures by constitutive family.}
\label{tab:energy_transport}
\begin{tabularx}{\textwidth}{@{}p{2.5cm}XX@{}}
\toprule
\textbf{Theory family} & \textbf{Energy/flux structure} & \textbf{Terms specific to the selected model} \\
\midrule
Classical thermoelasticity & $\mathcal{E}=\tfrac12\rho v_i v_i+\rho e$; exact Eulerian $F_i=q_i-\sigma_{ij}v_j+\mathcal{E}v_i$; reduced first-order $J_i=q_i-\sigma_{ij}v_j$ & Heat-flux constitutive law and boundary power \\
Generalized constitutive families & Model-specific stored energy and flux & Evolution equations, internal variables, boundary power, and dissipation terms \\
\bottomrule
\multicolumn{3}{@{}p{\textwidth}@{}}{\footnotesize No universal flux formula is implied for the generalized families.}\\
\end{tabularx}
\end{table}


\begin{figure}[htbp]
\centering
\includegraphics[width=0.9\textwidth]{fig14-energy-flux-channels.pdf}
\caption{Energy-balance conventions and open flux audit. The classical local thermoelastic flux is derived under the stated outward-positive convention; generalized multiphysics contributions and modal energy identities remain open pending model-specific derivations.}
\label{fig:energy_flux}
\end{figure}

\subsection{Notation summary for Section 10}

\begin{table}[htbp]
\centering
\small
\caption{Notation introduced in Section 10.}
\label{tab:notation-sec10}
\begin{tabularx}{\textwidth}{@{}p{1.8cm}X p{2.2cm}p{2.1cm}@{}}
\toprule
\textbf{Symbol} & \textbf{Meaning} & \textbf{Units} & \textbf{Equation} \\
\midrule
$\mathcal{E}$ & Total energy density & J\,m$^{-3}$ & \eqref{eq:10.1} \\
$\bm{J}$ & Energy flux vector & W\,m$^{-2}$ & \eqref{eq:10.2} \\
$\mathcal{P}$ & Mechanical power input & W\,m$^{-3}$ & \eqref{eq:10.2} \\
$J_i$ & Classical energy flux component & W\,m$^{-2}$ & \eqref{eq:10.3} \\
$\avg{J_i}_t$ & Time-averaged energy flux & W\,m$^{-2}$ & \eqref{eq:10.4} \\
$\mathcal{E}_\Sigma$ & Interfacial energy density & J\,m$^{-2}$ & \eqref{eq:10.5} \\
$\mathcal{P}_\Sigma$ & Interfacial power input & W\,m$^{-2}$ & \eqref{eq:10.5} \\
$R, T, A$ & Reflection, transmission, absorption coefficients & dimensionless & \eqref{eq:10.11} \\
\bottomrule
\end{tabularx}
\end{table}

\subsection{Open questions for Section 10}

\begin{enumerate}[label=\textbf{OQ10.\arabic*}]
\item What is the explicit time-domain energy flux and dissipation function for each of the families listed in Section~\ref{sec:10.4}, derived from that family's free energy, evolution laws, and boundary-power terms? \label{oq:10.1}
\item Is there a common formal structure under which the energy fluxes of the Lord--Shulman, phase-lag, Moore--Gibson--Thompson, and memory-dependent families become instances of a single template, so that a single derivation suffices for all of them? \label{oq:10.2}
\item Which observable would distinguish the energy partition predicted by two competing theories at the same interface, so that the difference could be measured rather than argued? \label{oq:10.3}
\item How should the entropy production of the conservative Green--Naghdi type~II structure be related to the entropy production of the dissipative branches, given that the former has no dissipation mechanism? \label{oq:10.4}
\item What normalisation of complex modes makes the energy ratios of Eq.~\eqref{eq:10.11} reproducible across independent implementations? \label{oq:10.5}
\item Can a dissipation potential be constructed for the memory-dependent kernel in the spirit of the classical one for the relaxation terms, and does its positivity restrict the admissible kernels to a smaller set than is currently assumed? \label{oq:10.6}
\item What is the relative importance of interface dissipation and bulk dissipation in the absorption coefficient of Eq.~\eqref{eq:10.11} for realistic bonded joints? \label{oq:10.7}
\end{enumerate}

% LocalWords:
